\section{Discussion and limitations}
\label{sec:discussion}

\subsection{When retaining a reference is worth the cost}
The fixed reference is an application constraint, not a universal performance objective. It can isolate scheduling changes during regression testing, permit exact replay from stored state, or support a component migration while its consumer remains unchanged. In such settings, a small execution improvement can be useful because the optimized path preserves the expected numerical interface. The 32B result shows that this option has measurable value beyond a strong implementation control. It does not establish a production return on maintenance effort, energy savings, or a general preference for the legacy backend.

A deployment free to change its numerical contract should evaluate that option directly. The modern extracted operator illustrates the possible benefit. A fair system decision would compare complete, appropriately qualified backends on the intended workload and include the cost of migration or compatibility. Our current evidence contains the legacy-compatible service comparison and the modern operator comparison, with different scopes; the missing full-backend comparison remains an open engineering decision rather than an assumed win for either choice.

\subsection{What transfers beyond this implementation}
The transferable object is the construction procedure: derive the physical and arithmetic constraints of a chosen reference, place the required cuts at the operations that depend on them, and check the consumer's outputs and successor state. It directs implementation and testing even when a backend uses different widths or cache layouts.

The implementation remains specialized. Its two projection widths and CUDA/Vulkan paths demonstrate instances of the construction, not automatic transfer to arbitrary quantizers, paged or sliding caches, mixed-device partitions, Flash attention, concurrent sequences, or stochastic sampling. Adding such a path requires deriving its numerical descriptor, discharging its state obligations and qualifying its actual dispatch. The conditional composition argument helps identify those obligations; it does not replace that work. Likewise, existing batching and packing techniques supply components, so the contribution should be judged by the combined execution organization and the evidence explaining its costs.

\subsection{Why the negative results matter}
The results expose three distinct constraints on performance. A local cut has little value when a request seldom crosses the relevant physical boundary. Even at a crossing, preserving a wider graph may leave the same arithmetic tiles or incur enough attention overhead to lose. Finally, a more efficient fixed operation can still make an application slower when it induces additional draft and target work. These mechanisms explain the small factorial cut effect, the multi-turn pilot and the wider-speculation reversal without requiring a universal threshold for useful speedup.

The evidence is also limited by its sampling. Several studies reuse development roots, and the Qwen scales reuse input identities. Confidence intervals preserve root clustering and describe the stated sessions; they do not measure cross-machine, cross-day or population-wide variation. Structural cells and the two-episode pilot are descriptive. Workload identities, orders, clocks and uncertainty types therefore remain explicit in the figures and supplement.

\subsection{Reproducibility and implementation maturity}
Public timing rows support independent statistical reconstruction, and the source maps identify the code and public diagnostics behind each result. Some native records are company-restricted, and some temporary intermediates were pruned. Public summaries cannot substitute for unavailable original tensors. The present reconstruction checks published statistics and implementations without rerunning native qualification. The final archival deposit and any permitted restricted-access procedure require author confirmation.

The current command-line implementation also has a narrower recovery boundary than a persistent inference service. Target-group errors poison the context, while other failures terminate the worker. A reusable service library would need transaction-wide invalidation and explicit admission for the unsupported graph cases. These requirements concern deployment maturity; the measured service timer remains the complete resident request path in the stated worker.
